\chapter{Efficient self-attention with spatial reduction}\label{ch:esa}

\begin{onething}
Keep the questions detailed but summarise the answers: queries stay at full resolution, while
keys and values are first shrunk by a strided convolution, so that every stage of the encoder
attends over the same $4\times4$ summary of 16 tokens instead of over every token.
\end{onething}

\Cref{ch:attention} built attention and ended on its cost: an $N\times N$ score matrix, over a
million entries at stage~1. This chapter derives that cost exactly, in multiply-adds and in bytes,
then shows the single modification SegFormer makes \citep{xie2021segformer}, borrowed from the
Pyramid Vision Transformer \citep{wang2021pyramid}: spatial-reduction attention. The code is
\file{app/foundation_models/components/efficient_self_attention.py}. Every number in the chapter
is printed by \file{research/textbook/tools/compute_13-15.py}.

\section{The cost of plain attention, counted}\label{sec:esa:cost}

\subsection{Multiply-adds}

We count work in \emph{multiply-adds}\index{multiply-add} (MACs): one multiplication followed by
one addition, the basic step of every dot product. A dot product of two length-$d$ vectors is
$d$ MACs.

Take one attention module on one patch with $N$ tokens of width $C$, split into $h$ heads of
width $d = C/h$ (\cref{sec:attention:heads}). Its steps cost:
\begin{align*}
  \text{each of } \mat{Q}, \mat{K}, \mat{V} &: N \text{ tokens} \times C\cdot C \text{ MACs} = NC^2 && \text{a } C\times C \text{ matrix per token}\\
  \mat{Q}\mat{K}\T \text{ per head} &: N\cdot N \text{ entries} \times d \text{ MACs} = N^2 d\\
  \mat{A}\mat{V} \text{ per head} &: N \text{ rows} \times d \text{ columns} \times N \text{ terms} = N^2 d\\
  \text{output projection} &: NC^2 .
\end{align*}
Add them, using $h\cdot d = C$ to sum the two per-head steps over the $h$ heads:
\begin{equation}\label{eq:esa:plaincost}
  \text{MACs}_\text{plain} = \underbrace{4NC^2}_{\text{four projections}} + \underbrace{2N^2C}_{\text{scores and blend}} .
\end{equation}
The first term grows in proportion to $N$; the second with $N^2$. At stage~1 the second wins
heavily.

\begin{worked}[label={we:esa:plaincost}]{The cost of plain attention at stage 1}
Stage 1 has $N = 1024$ tokens of width $C = 32$. Work each term out.
\begin{align*}
  C^2 &= 32\times32 = 1{,}024\\
  4NC^2 &= 4\times1024\times1{,}024 = 4\times1{,}048{,}576 = 4{,}194{,}304\\
  N^2 &= 1024\times1024 = 1{,}048{,}576\\
  2N^2C &= 2\times1{,}048{,}576\times32 = 2\times33{,}554{,}432 = 67{,}108{,}864\\
  \text{total} &= 4{,}194{,}304 + 67{,}108{,}864 = 71{,}303{,}168 .
\end{align*}
About 94\% of the work ($67.1/71.3$) is the score matrix and the blend.
\end{worked}

\subsection{Memory}

The score tensor has shape \shape{(B, h, N, N)} and the softmax weights another tensor of the same
shape. Training must keep the weights in memory until the backward pass (\cref{ch:trainloop})
needs them. Each float32 number is 4 bytes.

\begin{worked}[label={we:esa:plainmem}]{The memory of one stage-1 score tensor}
Per patch and per head, the scores are $1024\times1024 = 1{,}048{,}576$ numbers, and
\[
  1{,}048{,}576 \times 4 \text{ bytes} = 4{,}194{,}304 \text{ bytes} = 4 \text{ MiB}
\]
(a MiB is $2^{20} = 1{,}048{,}576$ bytes). The v0.2.0 pretraining configuration uses a batch of
$B = 128$ patches\tagM{configs/hyperspectral_segformer_exp_2.json} and stage 1 has $h = 2$ heads,
so one block's weights tensor would be
\begin{align*}
  128 \times 2 \times 4{,}194{,}304 \text{ bytes} &= 256 \times 4{,}194{,}304 \text{ bytes}\\
  &= 1{,}073{,}741{,}824 \text{ bytes} = 1{,}024 \text{ MiB} = 1 \text{ GiB},
\end{align*}
and stage 1 has two blocks, each storing scores and weights. Several gigabytes for one stage of
one batch, before counting anything else.
\end{worked}

\subsection{Why the obvious fixes fail}

\begin{enumerate}
  \item \emph{Use fewer, bigger tokens.} Stage 1 could cut the patch into $8\times8$ blocks
  instead of $4\times4$ ($N = 256$). But then every output token describes 64 pixels instead of
  16, and the fine detail that small anomalies live in is gone from the whole encoder.
  \item \emph{Attend only to nearby tokens.} Restrict each token to a window of neighbours. The
  cost falls, but a token in one corner can no longer see the other corner, and the context that
  makes a dark block surprising (``this is a desert'') may be outside the window.
  \item \emph{Average the keys and values over blocks.} Closer, but the averaging is fixed: it
  cannot learn that, inside a block, the edge pixels matter more than the middle ones.
\end{enumerate}
Spatial-reduction attention keeps every query (so every token still gets its own update), keeps
global reach (every query sees a summary of the whole patch), and makes the summarising
\emph{learned}.

\begin{analogy}
You have a precise question about a thousand-page report. Rather than reading all thousand pages
you read a sixteen-paragraph executive summary, one paragraph per chapter. Your question is as
specific as it was; only the material you consult has been condensed. And the summary was written
by someone who learned, over many reports, which parts of each chapter are worth keeping: that is
the learned convolution. Where the analogy stops: a person reads one summary once, while here
each of the 1{,}024 questions consults the same 16-paragraph summary, and every block writes a
fresh summary of its own.
\end{analogy}

\section{The strided convolution that summarises}\label{sec:esa:conv}

\subsection{Kernel equal to stride}

Recall from \cref{ch:shapes} that a 2-D convolution with kernel size $k$ and stride $s$ slides a
$k\times k$ window across the grid in steps of $s$, and that its output side length is
\begin{equation}\label{eq:esa:outsize}
  H_\text{out} = \left\lfloor \frac{H + 2p - k}{s} \right\rfloor + 1
\end{equation}
for padding $p$. The reduction uses $k = s = R$ and no padding ($p = 0$). Then
\begin{align*}
  H_\text{out} &= \left\lfloor \frac{H - R}{R} \right\rfloor + 1 && p = 0,\ k = s = R\\
  &= \left\lfloor \frac{H}{R} - 1 \right\rfloor + 1 && \text{split the fraction}\\
  &= \frac{H}{R} && \text{when } R \text{ divides } H .
\end{align*}
With kernel equal to stride the windows tile the grid exactly, with no overlap and no gap: each
output position summarises one $R\times R$ block of input tokens. (If $R$ did not divide $H$, the
floor would drop the last partial block: $H = 30$, $R = 8$ gives $\lfloor 22/8\rfloor + 1 = 3$, and
the last six rows would be ignored. For this model it always divides; see
\cref{sec:esa:sixteen}.)

\begin{definition}[Spatial-reduction convolution]
For token features $x_{c,i,j}$ (channel $c$, grid row $i$, column $j$) the reduced token at
summary position $(u, v)$ has channels
\begin{equation}\label{eq:esa:conv}
  y_{c',u,v} = b_{c'} + \sum_{c=1}^{C}\sum_{a=0}^{R-1}\sum_{b=0}^{R-1} W_{c',c,a,b}\; x_{c,\,Ru+a,\,Rv+b},
  \qquad c' = 1,\dots,C .
\end{equation}
\end{definition}\index{spatial reduction}\index{convolution!strided}
Every output channel looks at every input channel at every one of the $R^2$ positions of its
block. That is $C\cdot C\cdot R^2$ weights plus $C$ biases, a count we use in
\cref{sec:esa:params}.

\paragraph{Average pooling is one special case.} Choose $W_{c',c,a,b} = 1/R^2$ when $c' = c$
and 0 otherwise, and every bias $b_{c'} = 0$. Then \cref{eq:esa:conv} becomes
$y_{c,u,v} = \frac{1}{R^2}\sum_{a,b} x_{c,Ru+a,Rv+b}$, the plain average of the block, channel by
channel. The learned convolution can do this, or anything else linear: weight the centre more,
pick one corner, or mix channels.

\begin{worked}[label={we:esa:toyconv}]{Summarising a $4\times4$ grid with $R = 2$}
Take a grid of $4\times4$ tokens with a single channel ($C = 1$):
\[
  \begin{pmatrix} 1 & 2 & 0 & 1\\ 3 & 2 & 1 & 0\\ 0 & 0 & 4 & 4\\ 0 & 0 & 4 & 4 \end{pmatrix}.
\]
With $R = 2$ the output is $4/2 = 2$ by 2. The four blocks are the four $2\times2$ quarters.

\emph{Averaging kernel} $\begin{pmatrix}0.25 & 0.25\\ 0.25 & 0.25\end{pmatrix}$, bias 0:
\begin{align*}
  \text{top-left: } 0.25(1 + 2 + 3 + 2) &= 0.25\times8 = 2.0\\
  \text{top-right: } 0.25(0 + 1 + 1 + 0) &= 0.25\times2 = 0.5\\
  \text{bottom-left: } 0.25(0 + 0 + 0 + 0) &= 0.0\\
  \text{bottom-right: } 0.25(4 + 4 + 4 + 4) &= 0.25\times16 = 4.0
\end{align*}
\emph{``Pick the top-left'' kernel} $\begin{pmatrix}1 & 0\\ 0 & 0\end{pmatrix}$, bias 0: the outputs
are the top-left entries of the blocks, $1, 0, 0, 4$.

Two different summaries of the same grid. Training chooses the kernel, and with $C$ channels in
and out it chooses a different mix for each output channel. Sixteen tokens became four; with
$R = 8$ at stage 1, 1024 become 16.
\end{worked}

\begin{figure}[htbp]
\centering
\begin{tikzpicture}[cell/.style={minimum size=7mm, draw=white, line width=0.8pt, font=\small}]
  \foreach \r/\row in {0/{1,2,0,1},1/{3,2,1,0},2/{0,0,4,4},3/{0,0,4,4}}{
    \foreach \v [count=\c from 0] in \row {
      \pgfmathtruncatemacro{\p}{\v*18+6}
      \node[cell, fill=spat!\p] at (\c*0.7, -\r*0.7) {\v};
    }
  }
  \draw[spat!80!black, line width=1.2pt] (-0.35,0.35) rectangle (1.05,-1.05);
  \draw[spat!80!black, line width=1.2pt] (1.05,0.35) rectangle (2.45,-1.05);
  \draw[spat!80!black, line width=1.2pt] (-0.35,-1.05) rectangle (1.05,-2.45);
  \draw[spat!80!black, line width=1.2pt] (1.05,-1.05) rectangle (2.45,-2.45);
  \node[font=\small, align=center] at (1.05,0.8) {$4\times4$ tokens, $N = 16$};
  \draw[-{Latex}, thick] (2.8,-1.05) -- node[above, font=\small, align=center]{conv\\$k = s = 2$} (4.3,-1.05);
  \foreach \u/\v/\val in {0/0/2.0,1/0/0.5,0/1/0.0,1/1/4.0}{
    \pgfmathtruncatemacro{\p}{\val*18+6}
    \node[cell, minimum size=10mm, fill=spat!\p] at (5.0+\u*1.0, -0.55-\v*1.0) {\val};
  }
  \node[font=\small, align=center] at (5.5,0.8) {$2\times2$ summary};
  \node[font=\small, align=center] at (5.5,-2.5) {$N/R^2 = 4$ tokens};
  \draw[-{Latex}, thick] (6.8,-1.05) -- node[above, font=\small, align=center]{LN, then\\$\mat{W}_K$, $\mat{W}_V$} (8.6,-1.05);
  \node[font=\small, align=left, anchor=west] at (8.7,-1.05) {4 keys,\\4 values};
\end{tikzpicture}
\caption{The toy reduction of \cref{we:esa:toyconv} with the averaging kernel. Each output token
summarises one $R\times R$ block of the input (outlined); the queries are still made from all 16
input tokens.}
\label{fig:esa:toyconv}
\end{figure}

\subsection{LayerNorm after the reduction}

The summarised tokens pass through a LayerNorm\index{LayerNorm} before they are projected into
keys and values (\file{efficient_self_attention.py}, line 111). \Cref{ch:block} derives LayerNorm
in full; for now, it rescales each token so that its $C$ numbers have mean 0 and variance 1, then
applies a learned per-channel scale and shift. Here it keeps the summary tokens on the same scale
as the full-resolution tokens the queries come from, which matters because \cref{sec:attention:scale}
assumed unit-variance components when it chose the divisor $\sqrt d$.

\section{The data flow, with shapes at every step}\label{sec:esa:flow}

\begin{lstlisting}
# x: (B, N, C), N = H * W    (efficient_self_attention.py, lines 125-195)
Q = self.q(x)                                  # (B, N, C)
Q = Q.reshape(B, N, h, d).transpose(1, 2)      # (B, h, N, d): ALL tokens
r = x.transpose(1, 2)                          # (B, C, N)
r = r.reshape(B, C, H, W)                      # (B, C, H, W): a grid
r = self.reduction(r)                          # (B, C, H/R, W/R): conv
r = r.flatten(2).transpose(1, 2)               # (B, N/R^2, C): sequence
r = self.reduction_norm(r)                     # (B, N/R^2, C): LayerNorm
K = self.k(r).reshape(B, -1, h, d).transpose(1, 2)  # (B, h, N/R^2, d)
V = self.v(r).reshape(B, -1, h, d).transpose(1, 2)  # (B, h, N/R^2, d)
S = Q @ K.transpose(-2, -1) * d ** -0.5        # (B, h, N, N/R^2)
A = softmax(S, dim=-1)                         # (B, h, N, N/R^2)
O = A @ V                                      # (B, h, N, d)
O = O.transpose(1, 2).reshape(B, N, C)         # (B, N, C)
O = self.proj(O)                               # (B, N, C)
\end{lstlisting}
(This is a condensed transcription; the source spells each step out on its own line with the same
shapes in its comments.) Why the \code{transpose} before the \code{reshape}? Memory order: a
\shape{(B, N, C)} tensor stores token 0's $C$ numbers, then token 1's; a grid \shape{(B, C, H, W)}
must store channel 0's $H\cdot W$ numbers first. A reshape alone would scramble them
(\cref{ch:shapes}).

The key sentence for reading the scores: \emph{entry $(b, m, i, j)$ of $\mat{S}$ is how much
full-resolution token $i$ should attend, in head $m$, to summary token $j$, which stands for an
$R\times R$ block of the grid.} Fine-grained questions, spatially aggregated answers. The output
still has one vector per original token: nothing was lost on the query side.
\index{attention!efficient self-attention}\index{efficient self-attention}

\subsection{Every stage}

For a $128\times128$ patch, the stage widths, heads and ratios
$C = (32, 64, 160, 256)$, $h = (2, 2, 5, 8)$, $R = (8, 4, 2, 1)$ give these shapes
\par\noindent\tagM{research/experiments/indradhanu-finetune/FT04-final-joint/config.json}:

\begin{center}
\small
\setlength{\tabcolsep}{4pt}
\begin{tabular}{lllll}
\toprule
step & stage 1 & stage 2 & stage 3 & stage 4\\
 & $C{=}32, h{=}2, R{=}8$ & $C{=}64, h{=}2, R{=}4$ & $C{=}160, h{=}5, R{=}2$ & $C{=}256, h{=}8, R{=}1$\\
\midrule
input $\vect{x}$ & $(B, 1024, 32)$ & $(B, 256, 64)$ & $(B, 64, 160)$ & $(B, 16, 256)$\\
$\mat{Q}$ split & $(B, 2, 1024, 16)$ & $(B, 2, 256, 32)$ & $(B, 5, 64, 32)$ & $(B, 8, 16, 32)$\\
grid & $(B, 32, 32, 32)$ & $(B, 64, 16, 16)$ & $(B, 160, 8, 8)$ & not built\\
after conv & $(B, 32, 4, 4)$ & $(B, 64, 4, 4)$ & $(B, 160, 4, 4)$ & no conv\\
sequence, LN & $(B, 16, 32)$ & $(B, 16, 64)$ & $(B, 16, 160)$ & $(B, 16, 256)$\\
$\mat{K}, \mat{V}$ split & $(B, 2, 16, 16)$ & $(B, 2, 16, 32)$ & $(B, 5, 16, 32)$ & $(B, 8, 16, 32)$\\
scores $\mat{S}$ & $(B, 2, 1024, 16)$ & $(B, 2, 256, 16)$ & $(B, 5, 64, 16)$ & $(B, 8, 16, 16)$\\
$\mat{A}\mat{V}$ & $(B, 2, 1024, 16)$ & $(B, 2, 256, 32)$ & $(B, 5, 64, 32)$ & $(B, 8, 16, 32)$\\
output & $(B, 1024, 32)$ & $(B, 256, 64)$ & $(B, 64, 160)$ & $(B, 16, 256)$\\
\bottomrule
\end{tabular}
\end{center}
At stage 4, $R = 1$ and the module has no reduction layer at all (the \code{if} at line 103 only
builds it for $R > 1$); keys and values come straight from the 16 tokens, which is ordinary full
attention. The script runs a random input through each stage's module and gets the output shapes
in the last row\tagC.

\begin{worked}[label={we:esa:toyattn}]{One query against the four summary tokens}
Continue \cref{we:esa:toyconv} with the averaging summary $(2.0, 0.5, 0.0, 4.0)$. Keep $C = 1$,
one head ($d = 1$, so the scale is $1/\sqrt1 = 1$), and identity projections. (We skip the
LayerNorm here: with a single channel it would map every token to the same number, which is
\cref{ex:esa:5}.) Take the query of the top-left token, whose value is 1.

\emph{Scores,} query times each key: $1\times2.0 = 2.0$, $1\times0.5 = 0.5$, $1\times0.0 = 0.0$,
$1\times4.0 = 4.0$.

\emph{Stable softmax.} Subtract the maximum 4.0: $(-2.0, -3.5, -4.0, 0.0)$. Exponentiate:
\begin{align*}
  e^{-2.0} = 0.1353,\quad e^{-3.5} = 0.0302,\quad e^{-4.0} = 0.0183,\quad e^{0} = 1.0000 .
\end{align*}
Add: $0.1353 + 0.0302 + 0.0183 + 1.0000 = 1.1838$. Divide:
\begin{align*}
  0.1353/1.1838 = 0.1143,\quad 0.0302/1.1838 = 0.0255,\\
  0.0183/1.1838 = 0.0155,\quad 1.0000/1.1838 = 0.8447 .
\end{align*}
\emph{Blend the values:}
\begin{align*}
  0.1143\times2.0 &= 0.2286\\
  0.0255\times0.5 &= 0.0128\\
  0.0155\times0.0 &= 0.0000\\
  0.8447\times4.0 &= 3.3788\\
  \text{sum} &= 3.6202 .
\end{align*}
The top-left token is pulled strongly toward the bright bottom-right block, which it can see only
through its one summary token. The full layer does this for all 16 queries, so the score matrix is
$16\times4 = 64$ entries instead of $16\times16 = 256$: a factor of $R^2 = 4$.
\end{worked}

\coursefig[0.95\linewidth]{iw-13-esa}{Spatial reduction at stage 1, as the course draws it.
Queries keep all 1{,}024 tokens; keys and values are first shrunk by a strided convolution
(kernel = stride = $R = 8$) to a $4\times4$ grid of 16 summary tokens. The score matrix is
$1024\times16$ instead of $1024\times1024$, 64 times smaller (the hatched square is not drawn to
scale). With $R = 8, 4, 2, 1$ every stage attends over the same $4\times4$
summary.}{fig:esa:course}

\section{The 16-token summary}\label{sec:esa:sixteen}

\begin{keypoint}
The ratios $R = (8, 4, 2, 1)$ were chosen so that every stage, however fine its own grid, attends
over the same $4\times4$ summary of the whole patch.
\end{keypoint}

The token grid halves in side at each stage ($128/4 = 32$, then 16, 8, 4; \cref{ch:encoder}) and
the ratio halves too, so the summary side never changes:
\begin{align*}
  \text{stage 1: } 32/8 &= 4 && 4\times4 = 16 \text{ summary tokens}\\
  \text{stage 2: } 16/4 &= 4 && 4\times4 = 16\\
  \text{stage 3: } \phantom{1}8/2 &= 4 && 4\times4 = 16\\
  \text{stage 4: } \phantom{1}4/1 &= 4 && 4\times4 = 16 .
\end{align*}
In symbols: the stage-$s$ grid side is $128/(4\cdot2^{s-1})$ and $R_s = 8/2^{s-1}$, so the summary
side is $\frac{128}{4\cdot2^{s-1}}\cdot\frac{2^{s-1}}{8} = \frac{128}{32} = 4$, independent of $s$.
For a patch of side $P$ it would be $P/32$: a $256\times256$ patch would give an $8\times8$ summary
of 64 tokens at every stage.

\begin{center}
\begin{tabular}{cccccrrr}
\toprule
stage & grid & $N$ & $R$ & summary & $N/R^2$ & score entries & plain would be\\
\midrule
1 & $32\times32$ & 1024 & 8 & $4\times4$ & 16 & $1024\times16 = 16{,}384$ & 1{,}048{,}576\\
2 & $16\times16$ & 256 & 4 & $4\times4$ & 16 & $256\times16 = 4{,}096$ & 65{,}536\\
3 & $8\times8$ & 64 & 2 & $4\times4$ & 16 & $64\times16 = 1{,}024$ & 4{,}096\\
4 & $4\times4$ & 16 & 1 & $4\times4$ & 16 & $16\times16 = 256$ & 256\\
\bottomrule
\end{tabular}
\end{center}
The score matrix shrinks by exactly $R^2$: $64\times$ at stage 1, $16\times$ at stage 2, $4\times$
at stage 3, and not at all at stage 4, where 16 tokens is already the target.

\begin{figure}[htbp]
\centering
\begin{tikzpicture}
\begin{axis}[ybar, width=0.85\linewidth, height=6cm, ymode=log, log origin=infty,
  ymin=100, ymax=3e6, bar width=14pt,
  symbolic x coords={stage 1,stage 2,stage 3,stage 4}, xtick=data,
  ylabel={score entries per head per patch}, legend style={at={(0.98,0.97)},anchor=north east,font=\small},
  nodes near coords, every node near coord/.append style={font=\scriptsize, rotate=90, anchor=west},
  enlarge x limits=0.18]
\addplot[fill=maskgrey!60, draw=maskgrey] coordinates {(stage 1,1048576) (stage 2,65536) (stage 3,4096) (stage 4,256)};
\addplot[fill=attn!70, draw=attn] coordinates {(stage 1,16384) (stage 2,4096) (stage 3,1024) (stage 4,256)};
\legend{plain $N\times N$, reduced $N\times N/R^2$}
\end{axis}
\end{tikzpicture}
\caption{Score entries per head per patch, plain against spatially reduced, on a logarithmic axis
(each gridline is ten times the one below). The reduction factor $R^2$ is 64, 16, 4 and 1.}
\label{fig:esa:scorebars}
\end{figure}

\section{What the reduction saves, in work and memory}\label{sec:esa:savings}

\subsection{Multiply-adds with the reduction}

Write $N' = N/R^2$ for the number of summary tokens. Redo \cref{eq:esa:plaincost} step by step:
\begin{align*}
  \mat{Q} &: NC^2 && \text{all } N \text{ tokens, unchanged}\\
  \text{reduction conv} &: N' \text{ outputs} \times C \text{ channels} \times C R^2 \text{ MACs}\\
  &\phantom{:} = N'R^2C^2 = NC^2 && \text{since } N'R^2 = N\\
  \mat{K}, \mat{V} &: 2N'C^2 && \text{only } N' \text{ summary tokens}\\
  \mat{Q}\mat{K}\T,\ \mat{A}\mat{V} &: 2NN'C && N' \text{ keys instead of } N\\
  \text{output projection} &: NC^2 .
\end{align*}
So
\begin{equation}\label{eq:esa:redcost}
  \text{MACs}_\text{reduced} = 3NC^2 + 2N'C^2 + 2NN'C \qquad (R > 1).
\end{equation}
A curiosity: the reduction convolution costs exactly $NC^2$, the same as one projection, at every
stage. The LayerNorm adds a few operations per number and is left out.

\begin{worked}[label={we:esa:redcost}]{The cost of reduced attention at stage 1}
$N = 1024$, $N' = 16$, $C = 32$, $C^2 = 1{,}024$.
\begin{align*}
  3NC^2 &= 3\times1024\times1{,}024 = 3\times1{,}048{,}576 = 3{,}145{,}728\\
  2N'C^2 &= 2\times16\times1{,}024 = 32{,}768\\
  2NN'C &= 2\times1024\times16\times32 = 2\times524{,}288 = 1{,}048{,}576\\
  \text{total} &= 3{,}145{,}728 + 32{,}768 + 1{,}048{,}576 = 4{,}227{,}072 .
\end{align*}
Against the plain 71{,}303{,}168 of \cref{we:esa:plaincost}:
$71{,}303{,}168 / 4{,}227{,}072 = 16.9$ times less work. The score-and-blend part alone fell from
67{,}108{,}864 to 1{,}048{,}576, the full factor of 64; the projections, which do not shrink,
now dominate.
\end{worked}

The script evaluates both formulas at every stage\tagC:
\begin{center}
\begin{tabular}{crrrr}
\toprule
stage & plain MACs & reduced MACs & of which conv & ratio\\
\midrule
1 & 71{,}303{,}168 & 4{,}227{,}072 & 1{,}048{,}576 & 16.9\\
2 & 12{,}582{,}912 & 3{,}801{,}088 & 1{,}048{,}576 & 3.3\\
3 & 7{,}864{,}320 & 6{,}062{,}080 & 1{,}638{,}400 & 1.3\\
4 & 4{,}325{,}376 & 4{,}325{,}376 & 0 & 1.0\\
\bottomrule
\end{tabular}
\end{center}
After the reduction the four stages cost roughly the same. Without it, stage 1 alone would cost
nearly three times the other three stages together (71.3 million MACs against
$12.6 + 7.9 + 4.3 = 24.8$ million, all plain).

\subsection{Memory with the reduction}

\begin{worked}[label={we:esa:redmem}]{The stage-1 score tensor, reduced}
Per patch and head: $1024\times16 = 16{,}384$ numbers, times 4 bytes $= 65{,}536$ bytes
$= 64$ KiB, against 4 MiB plain. For the batch of 128 with 2 heads:
\begin{align*}
  128\times2\times65{,}536 \text{ bytes} = 256\times65{,}536 \text{ bytes} = 16{,}777{,}216 \text{ bytes} = 16 \text{ MiB},
\end{align*}
against 1{,}024 MiB plain. The same factor of 64.
\end{worked}

\section{What the reduction costs in parameters}\label{sec:esa:params}

The convolution of \cref{eq:esa:conv} has one weight for every (output channel, input channel,
kernel row, kernel column) and one bias per output channel; the LayerNorm has a scale and a shift
per channel:
\begin{equation}\label{eq:esa:convparams}
  \text{reduction conv} = C\cdot C\cdot R\cdot R + C,\qquad \text{reduction LN} = 2C .
\end{equation}
Add the four projections of \cref{eq:attention:params}:
\begin{equation}\label{eq:esa:params}
  \text{ESA}(C, R) = \underbrace{4C^2 + 4C}_{\text{q, k, v, proj}} + \underbrace{C^2R^2 + C + 2C}_{\text{only if } R > 1} .
\end{equation}

\begin{worked}[label={we:esa:params1}]{The stage-1 attention module, parameter by parameter}
$C = 32$, $R = 8$.
\begin{align*}
  \text{one projection: } 32\times32 + 32 &= 1{,}024 + 32 = 1{,}056\\
  \text{four projections: } 4\times1{,}056 &= 4{,}224\\
  \text{reduction conv: } 32\times32\times8\times8 + 32 &= 1{,}024\times64 + 32 = 65{,}536 + 32 = 65{,}568\\
  \text{reduction LN: } 2\times32 &= 64\\
  \text{total: } 4{,}224 + 65{,}568 + 64 &= 69{,}856 .
\end{align*}
\emph{Check.} Building \code{EfficientSelfAttention(32, 2, 8)} the script reports 69{,}856, with
\code{q}, \code{k}, \code{v}, \code{proj} at 1{,}056 each, \code{reduction} at 65{,}568 and
\code{reduction_norm} at 64\tagC.
\end{worked}

The same arithmetic at every stage, checked against the built modules\tagC:
\begin{center}
\begin{tabular}{crrrrrr}
\toprule
stage & $C$ & $R$ & projections & conv $C^2R^2 + C$ & LN $2C$ & ESA total\\
\midrule
1 & 32 & 8 & 4{,}224 & $65{,}536 + 32 = 65{,}568$ & 64 & 69{,}856\\
2 & 64 & 4 & 16{,}640 & $65{,}536 + 64 = 65{,}600$ & 128 & 82{,}368\\
3 & 160 & 2 & 103{,}040 & $102{,}400 + 160 = 102{,}560$ & 320 & 205{,}920\\
4 & 256 & 1 & 263{,}168 & none & none & 263{,}168\\
\bottomrule
\end{tabular}
\end{center}

\paragraph{The surprise.} At stage 1 the reduction convolution is $65{,}568/69{,}856 = 93.9\%$ of
the attention module, about fifteen times all four projections together
($65{,}568/4{,}224 = 15.5$). The reason is in \cref{eq:esa:convparams}: $C^2R^2 = (CR)^2$. At
stage 1, $CR = 32\times8 = 256$; at stage 2, $CR = 64\times4 = 256$ again, which is why the two
convolutions are almost the same size ($256^2 = 65{,}536$ weights each). Stage 1 has the narrowest
tokens and yet one of the larger attention modules. The 64-fold saving in compute is bought with
parameters; it is a trade, not a free lunch.

\begin{drift}
The widths, heads and ratios above are those of the FT04 configuration
(\file{research/experiments/indradhanu-finetune/FT04-final-joint/config.json}, read 2026-10-29).
Indradhanu Heavy triples the widths to $(96, 192, 480, 768)$ with heads $(3, 6, 10, 12)$
\tagM{indradhanu workings/68-indradhanu-heavy-pretraining.md}; its stage-1 reduction convolution is
$96\times96\times64 + 96 = 589{,}920$ parameters\tagC, nine times FT04's (\cref{ch:heavy}).
\end{drift}

\section{The fallback: when the grid is incomplete}\label{sec:esa:fallback}

The reduction rebuilds a full 2-D grid from the token sequence, which requires $N = H\cdot W$
exactly. During masked training (\cref{ch:masking}) stage 1 removes the tokens it will be asked to
predict before the blocks run. With the mask ratio 0.65
\par\noindent\tagM{research/experiments/indradhanu-finetune/FT04-final-joint/config.json},
\par\noindent about
\[
  1024\times(1 - 0.65) = 1024\times0.35 = 358.4 \approx 358
\]
tokens survive, and 358 tokens cannot be arranged into a $32\times32$ grid. The code checks first
(\file{efficient_self_attention.py}, line 144):
\begin{lstlisting}
can_reduce = (self.reduction_ratio > 1) and (N == H * W)
if can_reduce:
    ...                   # grid, conv, flatten, LayerNorm: K/V over N / R^2 tokens
else:
    reduced_x = x         # (B, N, C): K/V over all surviving tokens
\end{lstlisting}
So \emph{whenever tokens have been removed, stage 1 runs plain full attention over the
survivors.}\index{efficient self-attention!fallback}

\begin{worked}[label={we:esa:fallback}]{Is the fallback affordable?}
Per head and patch, with 358 survivors:
\[
  358\times358 = 128{,}164 \text{ score entries}.
\]
Compare: plain attention over all 1024 tokens, 1{,}048{,}576 (the fallback is
$128{,}164/1{,}048{,}576 = 12.2\%$ of it, about one eighth); reduced attention over all 1024 tokens,
16{,}384 (the fallback is $128{,}164/16{,}384 = 7.8$ times more). Removing 65\% of the tokens
shrinks $N^2$ by $0.35^2 = 0.1225$, which is why it is affordable.
\end{worked}

\begin{sourcenote}[: the fallback comment uses an old mask ratio]
The comment above the check (\file{efficient_self_attention.py}, lines 142--143) says ``With
\textasciitilde512 tokens after 50\% masking, full attention costs 512*512 = 262K''. That was the
mask ratio when the code was written. The FT04 configuration uses 0.65, giving about 358 survivors
and $358^2 = 128{,}164$ entries per head. The reasoning in the comment is right; its numbers are
dated. The course (IW13) already notes this.
\end{sourcenote}

Two consequences follow from reading the code, and both are worth knowing.
\begin{enumerate}
  \item \emph{The model is not the same machine with and without a mask.} A forward pass with a
  \code{keep_mask} at stage 1 uses full attention there; one without uses the reduction. This is
  intended and documented, but it means masked and unmasked outputs are produced by slightly
  different computations (\cref{ch:inference} shows how inference uses both).
  \item \emph{In a masked pass the stage-1 reduction layers are idle.} The \code{else} branch never
  touches \code{self.reduction} or \code{self.reduction_norm}, so a masked forward pass gives them
  no gradient. Their 65{,}632 parameters ($65{,}568 + 64$) are trained only by passes in which all
  1024 stage-1 tokens are present. Stages 2--4 never remove tokens, so their reductions always run.
\end{enumerate}

\begin{confusions}
\confusion{Which of query, key and value stay at full resolution?}{Only the queries. Every token
must receive its own updated vector, so there must be one query per token. Keys and values are
what is consulted, and those are summarised.}
\confusion{Does the reduction lose information?}{On the answer side, yes: fine spatial detail
within each $R\times R$ block is blended. The bet, validated by the SegFormer authors across many
tasks, is that queries need to be precise while answers can be summarised. The next stage of the
encoder works at a coarser grain anyway.}
\confusion{Is the reduction the same as pooling?}{It is learned pooling. Average pooling is one
setting of its weights (\cref{we:esa:toyconv}); training is free to find a better one, and to mix
channels while it summarises.}
\confusion{Does $R = 1$ at stage 4 mean attention is switched off?}{No. $R = 1$ means no reduction,
which is ordinary full attention over the 16 tokens. The module has no reduction layer at that
stage.}
\confusion{Does the fallback make masked training less accurate?}{It makes stage 1 less compressed,
if anything: full attention over the survivors sees more detail than a summary would. The cost is
compute, and that is small because most tokens are gone.}
\confusion{Why does ESA have more parameters at stage 1 than its width suggests?}{Because the
reduction convolution has $(CR)^2$ weights and $CR = 256$ at stage 1. The projections grow as
$C^2$; the convolution grows as $C^2R^2$.}
\end{confusions}

\begin{checkyourself}
\item\label{ex:esa:1} Stage 2 has $N = 256$ tokens on a $16\times16$ grid, $C = 64$, $h = 2$,
$R = 4$. Write the shape after each step of the data flow for a batch of $B$ patches, from the
input to the output projection.
\item\label{ex:esa:2} Count the parameters of the stage-3 reduction convolution and LayerNorm by
hand ($C = 160$, $R = 2$), and then the whole stage-3 attention module. Compare with 205{,}920.
\item\label{ex:esa:3} For a $256\times256$ patch, give the token grid, the number of summary
tokens and the score-matrix size per head at each of the four stages. By what factor has stage 1's
score matrix grown compared with a $128\times128$ patch?
\item\label{ex:esa:4} Evaluate \cref{eq:esa:plaincost,eq:esa:redcost} by hand for stage 2
($N = 256$, $N' = 16$, $C = 64$) and check them against the table in \cref{sec:esa:savings}.
\item\label{ex:esa:5} Show that LayerNorm on a token with a single channel ($C = 1$) always returns
its learned shift $\beta$, whatever the input. Why does this not matter for the real model?
\item\label{ex:esa:6} With the older mask ratio 0.5, how many stage-1 tokens survive, how many
score entries does the fallback compute per head, and what fraction is that of plain attention over
all 1024 tokens?
\end{checkyourself}

\nextchapter{ch:block}
